% !TEX spellcheck = en_US
% !TEX spellcheck = LaTeX
\documentclass[letterpaper,10pt,english]{article}
\usepackage{%
amsfonts,%
amsmath,%
amssymb,%
amsthm,%
babel,%
bbm,%
%biblatex,%
caption,%
centernot,%
color,%
enumerate,%
%enumitem,%
epsfig,%
epstopdf,%
etex,%
fancybox,%
framed,%
fullpage,%
%geometry,%
graphicx,%
hyperref,%
latexsym,%
mathptmx,%
mathtools,%
multicol,%
paralist,%
pgf,%
pgfplots,%
pgfplotstable,%
pgfpages,%
proof,%
psfrag,%
%subfigure,%
tcolorbox,%
tfrupee,%
tikz,%
times,%
%ulem,%
url,%
xcolor,%
mathpazo,%
}

\definecolor{shadecolor}{gray}{.95}%{rgb}{1,0,0}
\usepackage[margin=1in,top=0.75in]{geometry}
\usepackage[mathscr]{eucal}
\usepgflibrary{shapes}
\usepgfplotslibrary{fillbetween}
\usetikzlibrary{%
arrows,%
backgrounds,%
chains,%
decorations.pathmorphing,% /pgf/decoration/random steps | erste Graphik
decorations.text,% 
matrix,%
positioning,% wg. " of "
fit,%
patterns,%
petri,%
plotmarks,%
scopes,%
shadows,%
shapes.misc,% wg. rounded rectangle
shapes.arrows,%
shapes.callouts,%
shapes%
}

%\pgfplotsset{compat=newest} %<------ Here
\pgfplotsset{compat=1.11} %<------ Or use this one

\theoremstyle{plain}
\newtheorem{thm}{Theorem}[section]
\newtheorem{lem}[thm]{Lemma}
\newtheorem{prop}[thm]{Proposition}
\newtheorem{cor}[thm]{Corollary}
\newtheorem{clm}[thm]{Claim}

\theoremstyle{definition}
\newtheorem{defn}[thm]{Definition}
\newtheorem{conj}[thm]{Conjecture}
\newtheorem{exmp}[thm]{Example}
\newtheorem{axiom}[thm]{Axiom}
\newtheorem{assum}[thm]{Assumption}
\newtheorem{exerc}[thm]{Exercise}
\newtheorem*{defin}{Definition}

\theoremstyle{remark}
\newtheorem{rem}{Remark}
\newtheorem{note}{Note}

\newcommand{\iid}{\emph{i.i.d.}\ }
\newcommand{\Cov}{\operatorname{Cov}}
%\newcommand{\det}{\operatorname{det}}
%\newcommand{\Real}{\mathbb{R}}
\newcommand{\tr}{\operatorname{tr}}
\newcommand{\Var}{\operatorname{Var}}
\newcommand{\cov}{\operatorname{cov}}

%\renewcommand{\proof}[1]{\begin{proof}#1\end{proof}}
\newcommand{\EQ}[1]{\begin{equation*}#1\end{equation*}}
\newcommand{\EQN}[1]{\begin{equation}#1\end{equation}}
\newcommand{\eq}[1]{\begin{align*}#1\end{align*}}
\newcommand{\meq}[2]{\begin{xalignat*}{#1}#2\end{xalignat*}}
\newcommand{\norm}[1]{\left\lVert#1\right\rVert}
\newcommand{\inner}[1]{\left\langle #1\right\rangle}
\newcommand{\abs}[1]{\left\lvert#1\right\rvert}
\newcommand{\expect}[1]{\mathbb{E}\left[{#1}\right]}
\newcommand{\prob}[1]{\mathbb{P}\left[{#1}\right]}
\newcommand{\given}{\; \big\vert \;} 
\newcommand{\set}[1]{\left\{#1\right\}} 
\newcommand{\indicator}[1]{1_{\set{#1}}} 
\newcommand{\Ind}[1]{\mathbbm{1}_{#1}} 
\newcommand{\SetIn}[1]{\mathbbm{1}_{\set{#1}}} 
\newcommand{\red}[1]{\textcolor{red}{#1}} 
\newcommand{\floor}[1]{\left\lfloor#1\right\rfloor}
\newcommand{\ceil}[1]{\left\lceil#1\right\rceil}


\newcommand{\C}{\mathbb{C}}
\newcommand{\D}{\mathbb{D}}
\newcommand{\E}{\mathbb{E}}
\newcommand{\F}{\mathbb{F}}
\newcommand{\N}{\mathbb{N}}
\renewcommand{\P}{\mathbb{P}}
\newcommand{\Q}{\mathbb{Q}}
\newcommand{\R}{\mathbb{R}}
\newcommand{\Z}{\mathbb{Z}}

\newcommand{\bA}{\mathbf{A}}
\newcommand{\bI}{\mathbf{I}}
\newcommand{\bU}{\mathbf{1}}
\newcommand{\bx}{\mathbf{x}}
\newcommand{\bX}{\mathbf{X}}
\newcommand{\bZ}{\mathbf{Z}}


\newcommand{\cA}{\mathcal{A}}
\newcommand{\cB}{\mathcal{B}}
\newcommand{\cC}{\mathcal{C}}
\newcommand{\cF}{\mathcal{F}}
\newcommand{\cG}{\mathcal{G}}
\newcommand{\cM}{\mathcal{M}}
\newcommand{\cN}{\mathcal{N}}
\newcommand{\cP}{\mathcal{P}}
\newcommand{\cT}{\mathcal{T}}
\newcommand{\cX}{\mathcal{X}}
\newcommand{\cY}{\mathcal{Y}}

\newcommand{\sA}{\mathscr{A}}
\newcommand{\sB}{\mathscr{B}}
\newcommand{\sC}{\mathscr{C}}
\newcommand{\sE}{\mathscr{E}}
\newcommand{\sF}{\mathscr{F}}
\newcommand{\sG}{\mathscr{G}}
\newcommand{\sH}{\mathscr{H}}
\newcommand{\sL}{\mathscr{L}}
\newcommand{\sS}{\mathscr{S}}
\newcommand{\sT}{\mathscr{T}}
\newcommand{\sX}{\mathscr{X}}
\newcommand{\sY}{\mathscr{Y}}
\newcommand{\sZ}{\mathscr{Z}}

\newcommand{\tS}{\tilde{S}}
% Debug
\newcommand{\todo}[1]{\begin{color}{blue}{{\bf~[TODO:~#1]}}\end{color}}

% a few handy macros

\renewcommand{\le}{\leqslant}
\renewcommand{\ge}{\geqslant}
\newcommand\matlab{{\sc matlab}}
\newcommand{\goto}{\rightarrow}
\newcommand{\bigo}{{\mathcal O}}
%\newcommand{\half}{\frac{1}{2}}
%\newcommand\implies{\quad\Longrightarrow\quad}
\newcommand\reals{{{\rm l} \kern -.15em {\rm R} }}
\newcommand\complex{{\raisebox{.043ex}{\rule{0.07em}{1.56ex}} \hskip -.35em {\rm C}}}


% macros for matrices/vectors:

% matrix environment for vectors or matrices where elements are centered
\newenvironment{mat}{\left[\begin{array}{ccccccccccccccc}}{\end{array}\right]}
\newcommand\bcm{\begin{mat}}
\newcommand\ecm{\end{mat}}

% matrix environment for vectors or matrices where elements are right justifvied
\newenvironment{rmat}{\left[\begin{array}{rrrrrrrrrrrrr}}{\end{array}\right]}
\newcommand\brm{\begin{rmat}}
\newcommand\erm{\end{rmat}}

% for left brace and a set of choices
%\newenvironment{choices}{\left\{ \begin{array}{ll}}{\end{array}\right.}
\newcommand\when{&\text{if~}}
\newcommand\otherwise{&\text{otherwise}}
% sample usage:
%\delta_{ij} = \begin{choices} 1 \when i=j, \\ 0 \otherwise \end{choices}


% for labeling and referencing equations:
\newcommand{\eql}{\begin{equation}\label}
\newcommand{\eqn}[1]{(\ref{#1})}
% can then do
%\eql{eqnlabel}
%...
%\end{equation}
% and refer to it as equation \eqn{eqnlabel}.
% some useful macros for finite difference methods:
\newcommand\unp{U^{n+1}}
\newcommand\unm{U^{n-1}}
% for chemical reactions:
\newcommand{\react}[1]{\stackrel{K_{#1}}{\rightarrow}}
\newcommand{\reactb}[2]{\stackrel{K_{#1}}{~\stackrel{\rightleftharpoons}
 {\scriptstyle K_{#2}}}~}
\makeatletter
\def\th@plain{%
\thm@notefont{}% same as heading font
\itshape % body font
}
\def\th@definition{%
\thm@notefont{}% same as heading font
\normalfont % body font
}
\makeatother
\date{}


\title{Lecture-27: Poisson process on the half-line}
\author{}

\begin{document}
\maketitle

%\section{Equivalent characterizations}
%A counting process $N$ has the \textbf{completely independence property},  
%if for any collection of finite disjoint and bounded sets $A_1, \dots, A_k \in \sB$, 
%\EQ{
%P\bigcap_{i=1}^k\set{N(A_i) = n_i} = \prod_{i=1}^kP\set{N(A_i) = n_i}.
%}
%\begin{thm} 
%\label{thm:void}
%Distribution of a simple point process is completely determined by void probabilities. 
%\end{thm}
%\begin{proof}
%We will show this by induction on the number of points in a bounded set $A \in \sB$. 
%We assume that $(A_k \in \sB: k \in \N)$ is a sequence of sufficiently small sets partitioning $A$ such that $N(A_i) \in \set{0,1}$ for all $i \in \N$. 
%%We can write 
%%\EQ{
%%P\set{N(A) = 1} = \sum_{k \in \N}P\set{N(A_k)=1, N(A\setminus A_k) = 0}. 
%%}
%%We observe the following set equality 
%%$
%%\set{N(A_k)=1, N(A\setminus A_k) = 0}\cup\set{N(A) = 0} = \set{N(A \setminus A_k) = 0},
%%$
%%and hence we can write 
%%\EQ{
%%P\set{N(A) = 1} = \sum_{k \in \N}\left(P\set{N(A \setminus A_k)=0}-P\set{N(A) = 0}\right). 
%%}
%We assume that $N(A) = n$ and by the induction hypothesis $P\set{N(B) = n-1}$ can be completely characterized by the void probabilities for all bounded sets $B \in \sB$. 
%Then, we can write
%\EQ{
%P\set{N(A) = n} = \sum_{k \in \N}P\set{N(A_k) = 1, N(A \setminus A_k) = n-1} = \sum_{k \in \N}\left(P\set{N(A \setminus A_k)=n-1}-P\set{N(A) = n-1}\right). 
%}
%\end{proof}
%\begin{thm}[Equivalences] 
%\label{thm:equiv}
%Following are equivalent for a simple counting process $N: \Omega \to {\Z_+}^\sB$. 
%\begin{enumerate}[i\_]
%\item Process $N$ is Poisson with locally finite intensity measure $\Lambda$. 
%\item For each bounded $A \in \sB$, we have $P\set{N(A) = 0} = e^{-\Lambda(A)}$. 
%\item %The intensity measure $\Lambda$ is bounded for bounded $A \in \sB$,  and 
%For each bounded $A \in \sB$,  the number of points $N(A)$ is a Poisson with parameter $\Lambda(A)$. 
%\item Process $N$ has the completely independence property, and $\E N(A) = \Lambda(A)$. 
%\end{enumerate}
%\end{thm}
%\begin{proof}
%We will show that $i\_ \implies ii\_ \implies iii\_ \implies iv\_ \implies i\_$. 
%\begin{enumerate}
%	\item[$i \implies ii\_$] It follows from the definition of Poisson point processes and definition of Poisson random variables.  
%	\item[$ii \implies iii\_$] From Theorem~\ref{thm:void}, we know that void probabilities determine the entire distribution. 
%	Further, we observe that 
%	\EQ{
%	\sum_{k\in \N}P\set{N(A) = k} = e^{-\Lambda(A)}\sum_{k\in \N}\frac{\Lambda(A)^k}{k!}.
%	}
%	\item[$iii \implies iv\_$] We will show this in two steps. 
%	\begin{itemize}
%	\item[Mean:] Since the distribution of random variable $N(A)$ is Poisson, it has mean $\E N(A) = \Lambda(A)$. 
%	\item[CIP: ] For disjoint and bounded $A_1, \dots, A_k \in \sB$ and $A = \cup_{i=1}^kA_i$, we have $N(A) = N(A_1) + \dots N(A_1)$. 
%	Taking expectations on both sides, and from the linearity of expectation, we get 
%	\EQ{
%	\Lambda(A) = \Lambda(A_1) + \dots + \Lambda(A_k).
%	} 
%	From the number of partitions $n_1 + \dots + n_k = n$, we can write 
%	\EQ{
%	P\set{N(A)=n} = \frac{1}{n!}\sum_{n_1 + \dots + n_k = n}\binom{n}{n_1, \dots, n_k}P\set{N(A_1) = n_1, \dots, N(A_k) = n_k}. 
%	}
%	Using the definition of Poisson distribution, we can write the left hand side of the above equation as 
%	\EQ{
%	P\set{N(A) = n} = e^{-\Lambda(A)} \frac{\Lambda(A)^n}{n!} = \prod_{i=1}^ke^{-\Lambda(A_i)} \frac{(\sum_{i=1}^k\Lambda(A_i))^n}{n!} =  \frac{1}{n!}\sum_{n_1 + \dots + n_ k =n}\binom{n}{n_1, \dots, n_k}\prod_{i=1}^ke^{-\Lambda(A_i)} \Lambda(A_i))^{n_i}. 
%	}
%	Equating each term in the summation, we get	
%	\EQ{
%	P\set{N(A_1) = n_1, \dots, N(A_k) = n_k} = \prod_{i=1}^kP\set{N(A_i) = n_i}. 
%	}
%	\end{itemize}
%
%	\item[$iv \implies i\_$] Due to complete independence property and since void probabilities describe the entire distribution, 
%	it suffices to show that $P\set{N(A) = 0} = e^{-\Lambda(A)}$ for all bounded $A \in \sB$. 
%	For disjoint and bounded $A_1, \dots, A_k \in \sB$ and $A = \cup_{i=1}^kA_i$, 
%	we have 
%	\meq{2}{
%	&\Lambda(A) = \sum_{i=1}^k\Lambda(A_i),&\text{ and }&-\ln P\set{N(A) = 0} = -\sum_{i=1}^k \ln P\set{N(A_i) = 0}.
%	} 
%	This implies that $-\ln P\set{N(A) = 0} = \Lambda(A)$, and the result follows. 
%\end{enumerate}
%\end{proof}
%
%%\red{
%%\begin{thm} 
%%Following are equivalent for a simple counting process $N$. 
%%\begin{enumerate}[i\_]
%%\item Process $N$ is Poisson with intensity measure $\Lambda$. 
%%\item Process $N$ has the completely independence property, and $\E N(A) = \Lambda(A)$. 
%%\item For any bounded $A \in \sB$, the intensity measure $\Lambda$ is bounded and $N(A)$ is Poisson with parameter $\Lambda(A)$. 
%%\end{enumerate}
%%\end{thm}
%%\begin{proof} 
%%We will show that $i\_ \implies ii\_ \implies iii\_ \implies I\_$. 
%%\begin{enumerate}[i\_]
%%\item[$i\_ \implies ii\_$] Let the process $N$ be Poisson with intensity measure $\Lambda$. 
%%Then the complete independence property follows from the definition, 
%%and since $N(A)$ is a Poisson distributed random variable with mean $\Lambda(A)$. 
%%\item[$iI\_ \implies iii\_$] Consider a process $N$ with the completely independence property, and $\E N(A) = \Lambda(A)$. 
%%Let $A \in \sB$ be bounded, then it follows that $\Lambda(A)$ is bounded. 
%%From the fact that $\Lambda(A) = \sum_{i=1}^k\Lambda(A_i)$ for any partition $(A_1, \dots, A_k)$ of the set $A$, 
%%we can write $\set{N(A) = 0} = \cap_{i=1}^k\set{N(A_i) = 0}$. 
%%Therefore, from the complete independence property, we have 
%%\EQ{
%%P\set{N(A) = 0} = \prod_{i=1}^kP\set{N(A_i) = 0}. 
%%}
%%This is the semi-group property, from which it follows that 
%%\meq{2}{
%%&P\set{N(A) = 0} = e^{-\Lambda(A)},&&\sum_{k\in \N}P\set{N(A) = k} = e^{-\Lambda(A)}\sum_{k\in \N}\frac{\Lambda(A)^k}{k!}.
%%}
%%From which it follows that $N(A)$ is Poisson with parameter $\Lambda(A)$. 
%%\item[$iii\_ \implies i\_$] Let $N$ be a simple counting process and $\Lambda$ an intensity measure such that for any $A \in \sB$, 
%%$\Lambda(A)$ is bounded and $N(A)$ is Poisson with parameter $\Lambda(A)$. 
%%This implies that for any partition $(A_1, \dots, A_k)$ of the set $A_i$, we have 
%%\EQ{
%%P\set{N(A) = 0} = \prod_{i=1}^kP\set{N(A_i) = 0}. 
%%}
%%Hence, we can show that the counting process $N$ has the complete independence property. 
%%\end{enumerate}
%%\end{proof}
%%}
%%%%%%%%%%%%%%%%%%%%%%%%%%%%%%%%%%%%%%%%%%%%%%%%%%%%%%%%%%%%%%%%%%%%%%%
\section{Simple point processes on the half-line} 
%%%%%%%%%%%%%%%%%%%%%%%%%%%%%%%%%%%%%%%%%%%%%%%%%%%%%%%%%%%%%%%%%%%%%%%
A stochastic process defined on the half-line $N: \Omega \to \Z_+^{\R_+}$ is a \textbf{counting process} if
\begin{enumerate}
  \item $N_0 = 0$, and 
  \item for each $\omega \in \Omega$, the sample path $N(\omega): \R_+ \to \Z_+$ is non-decreasing, integer valued, and right continuous function of time $t \in \R_+$.% and at points of discontinuity $(N_t- N(t^{-}))\le 1,~~\forall \omega \in \Omega$. 
\end{enumerate}
Each discontinuity of the sample path of the counting process can be thought of as a jump of the process, as shown in Figure~\ref{Fig:Poisson}.  
A simple counting process has the unit jump size almost surely. 
%This is equivalent to a point process not having an atom at any point on the line. 
General point processes in higher dimension don't have any inter-arrival time interpretation. 
%\end{rem}
\begin{figure}[hhhh]
\center
\scalebox{0.8}{\input{../../Figures/Poisson}}
\caption{Sample path of a simple counting process.}
\label{Fig:Poisson}
\end{figure}
%\begin{lem} 
%A counting process $N: \Omega \to \Z_+^{\R_+}$ has finitely many jumps in a finite interval $(0,t]$ almost surely. 
%\end{lem}
\begin{defn} 
The points of discontinuity are also called the \textbf{arrival instants} of the counting process $N$. 
The \textbf{$n$th arrival instant} is a random variable denoted $\tS_n: \Omega \to \R_+$, 
defined inductively as 
\meq{2}{
&\tS_0 \triangleq 0, &&\tS_n \triangleq \inf\set{t \ge 0: N_t \ge n},\quad n \in \N.
}
\end{defn}
\begin{defn} 
The \textbf{inter arrival time} between $(n-1)${th} and $n${th} arrival is denoted by $X_n$ and written as
%\EQ{
$X_n \triangleq \tS_n - \tS_{n-1}$. 
%}
\end{defn}
\begin{rem} 
For a simple point process, we have
%\EQ{
$P\set{X_{n} = 0} = P\set{X_n\le 0} = 0.$
%}
\end{rem}
\begin{lem} 
Simple counting process $N:\Omega\to\Z_+^{\R_+}$ and arrival process $\tS:\Omega\to\R_+^\N$ are inverse processes, i.e. 
\EQ{
\set{\tS_n \le t} = \set{N_t \ge n}.
}
\end{lem}
\begin{proof} 
Let $\omega \in \set{\tS_n \le t}$, then $N_{\tS_n} = n$ by definition. 
%from right continuity of the counting process. 
Since $N$ is a non-decreasing process, we have $N_t \ge N_{\tS_n} = n$. 
%It follows that $\set{S_n \le t} \subseteq \set{N_t \ge n}$.
Conversely, let $\omega \in \set{N_t \ge n}$, then it follows from definition that $\tS_n \le t$.
\end{proof}
\begin{cor} 
For arrival instants $\tS:\Omega\to\R_+^\N$ associated with a counting process $N:\Omega\to\Z_+^{\R_+}$ 
we have 
$
\set{\tS_n \le t, \tS_{n+1} > t} = \set{N_t = n} 
$
for all $n \in \Z_+$ and $t \in \R_+$. 
\end{cor}
\begin{proof}
It is easy to see that 
%\EQ{
$\set{\tS_{n+1} > t } = \set{\tS_{n+1} \le t}^c = \set{N_t \ge n+1}^c = \set{N_t < n+1}$.
%}
Hence, %the result follows by writing 
\EQ{
\set{N_t = n} = \set{N_t \ge n, N_t < n+1} = \set{\tS_n \le t, \tS_{n+1} > t}.
}
\end{proof}
\begin{lem}
Let $F_n(x)$ be the distribution function for $S_n$, then 
%\EQ{
$P_n(t) \triangleq P\set{N_t = n} = F_{n}(t)-F_{n+1}(t)$. 
%}
\end{lem}
\begin{proof} It suffices to observe that following is a union of disjoint events,
\EQ{
\set{\tS_n \le t } = \set{\tS_n \le t, \tS_{n+1} > t} \cup \set{\tS_{n} \le t, \tS_{n+1} \le t}.
}
\end{proof}


%\begin{cor}[Poisson process on the half-line]
%A random process $N: \Omega \to \Z_+^{\R_+}$ indexed by time $t \in \Z_+$ is a Poisson process with intensity measure $\Lambda$ iff
%\begin{enumerate}[(a)]
%\item Starting with $N(0) = 0$, the process $N_t$ takes a non-negative integer value for all $t \in \R_+$;
%\item the increment $N(t + s) - N_t$ is surely nonnegative for any $s \in \R_+$;
%\item the increments $N(t_1), N(t_2) - N(t_1),  \dots , N(t_n) - N(t_{n-1})$ are \textit{independent} for any 
%$0 < t_1 < t_2 < \dots < t_{n-1} < t_n$; 
%\item the increment $N(t+s) - N_t$ is distributed as Poisson random variable with parameter $\Lambda((t, t+s])$. 
%\end{enumerate}
%The Poisson process is homogeneous with intensity $\lambda$, iff in addition to conditions $(a), (b), (c)$, 
%the distribution of the increment $N(t + s) - N_t$ depends on the value $s \in \R_+$ but is independent of $t \in \R_+$. 
%That, is the increments are \textit{stationary}. 
%\end{cor}
%\begin{proof} 
%We have already seen that definition of Poisson processes implies all four conditions. 
%Conditions $(a)$ and $(b)$ imply that $N$ is a simple counting process on the half-line, 
%condition $(c)$ is the complete independence property of the point process, 
%and condition $(d)$ provides the intensity measure. 
%The result follows from the equivalence $iv\_$ in Theorem~\ref{thm:equiv}. 
%\end{proof}


%%%%%%%%%%%%%%%%%%%%%%%%%%%%%%%%%%%%%%%%%%%%%%%%%%%%%%%%%%%%%%%%%%%%%%%
\section{IID exponential inter-arrival times characterization}
%%%%%%%%%%%%%%%%%%%%%%%%%%%%%%%%%%%%%%%%%%%%%%%%%%%%%%%%%%%%%%%%%%%%%%%

\begin{prop} 
The counting process $N:\Omega\to\Z_+^{\R_+}$ 
associated with a simple Poisson point process $S:\Omega\to\R_+^\N$ 
is Markov. 
\end{prop}
\begin{proof} 
We define the event space $\sF_t \triangleq \sigma(N_s: s \le t)$ as the history of the process until time $t \in \R_+$. 
Then, from the independent increment property of Poisson processes, we have for any historical event $H_s \in \sF_s$ 
\EQ{
P(\set{N_t = n}\given H_s\cap\set{N_s = k})= P(\set{N_t-N_s = n-k} \given H_s\cap \set{N_s = k}) = P(\set{N_t = n}\given \set{N_s = k}). 
}
The transition probability matrix is $P(s,t)$ with its $(k,n)$th entry given by $e^{-\Lambda(s,t]}\frac{(\Lambda(s,t])^{n-k}}{(n-k)!}$. 
%The equality follows from the independence of increments. 
\end{proof}
\begin{rem}
A Markov process $X:\Omega\to\sX^\R$ is time homogeneous if the transition matrix $P(s,t) = P(t-s)$ for all $t \ge s$. 
Thus the counting process for a homogeneous Poisson point process is time homogeneous Markov process, as the transition probability matrix $P(s,t) = P(t-s)$ with its $(k,n)$th entry given by $e^{-\lambda(t-s)}\frac{(\lambda(t-s))^{n-k}}{(n-k)!}$. 
\end{rem}
\begin{thm} 
The counting process $N:\Omega\to\Z_+^{\R_+}$ 
associated with a simple Poisson point process $S:\Omega\to\R_+^\N$ 
is strongly Markov. 
\end{thm}
\begin{prop}
A simple counting process $N:\Omega\to\Z_+^{\R_+}$ is associated with a \textbf{homogeneous Poisson process} with a constant intensity density $\lambda$, 
iff the inter-arrival time sequence $X:\Omega\to\R_+^\N$ are \iid random variables with an exponential distribution of rate $\lambda$. 
%That is, it has a distribution function $F: \R_+ \to [0,1]$, such that $F(x) = 1 - e^{-\lambda x}$ for all $x \in \R_+$. 
\end{prop}
\begin{proof}
Let $N_t$ be a counting process associated with a homogeneous Poisson point process on half-line with constant intensity density $\lambda$. 
From equivalence $iii\_$ in Theorem~\ref{thm:equiv}, we obtain for any positive integer $t$, 
\EQ{
P\set{X_1 > t} = P\set{N_t = 0} = e^{-\lambda t}. 
}
It suffices to show that inter-arrivals time sequence $X:\Omega\to\R_+^\N$ is \iid.  
We can show that $N$ is Markov process with strong Markov property. 
Since the sequence of ordered points $\tS:\Omega\to\R_+^\N$ is a sequence of stopping times for the counting process, 
it follows from the strong Markov property of this process that $(N_{\tS_n+t} - N_{\tS_n}: t \ge 0)$ is independent of $\sigma(N_s: s \le \tS_n)$ and hence of $\tS_n$ and $N_{\tS_n}$. 
Further, we see that
\EQ{
X_{n+1} = \inf\set{t > 0: N_{\tS_n+t} - N_{\tS_n}  = 1}. 
}
It follows that $X:\Omega\to\R_+^\N$ is an independent sequence. 
For homogeneous Poisson point process, we have $N_{\tS_n+t}-N_{\tS_n} = N_t$ in distribution, 
and hence $X_{n+1}$ has same distribution as $X_1$ for each $n \in \N$. 

%We can show that $X_n = S_n - S_{n-1}$ is independent of $(N(S_n+t) - N(S_n): t \ge 0)$
For the given \iid inter-arrival time sequence $X:\Omega\to\R_+^\N$ distributed exponentially with rate $\lambda$, we define the $n$th arrival instant $\tS_n \triangleq \sum_{i=1}^nX_i$ for each $n \in \N$, and the number of arrivals in time duration $(0,t]$ as $N_t\triangleq \sum_{n\in\N}\SetIn{\tS_n\le t}$ for all $t\in \R_+$. 
It follows that $N_t$ is path wise non-decreasing, integer-valued, right continuous, and simple since $P\set{X_1 \le 0} = 0$. 
Therefore, $N$ is a simple counting process such that 
\EQ{
P\set{N_t = 0} = P\set{X_1 > t} = e^{-\lambda t}. 
}
It follows that the void probabilities are exponential and hence the random variable $N_t$ is Poisson with parameter $\lambda t$ for all $t \in \R_+$. 
Hence, $N$ is a counting process associated with a homogeneous Poisson process with the constant intensity density $\lambda$ from the equivalence $ii\_$ in Theorem~\ref{thm:equiv}. 
\end{proof}
%\begin{defn} 
%A simple counting process $(N_t: t\ge 0) $ is called a homogeneous \textbf{Poisson process} with a finite positive rate $\lambda$, 
%if the inter-arrival times $(X_{n}: n \in \N)$ are \iid random variables with an exponential distribution of rate $\lambda$. 
%That is, it has a distribution function $F: \R_+ \to [0,1]$, such that $F(x) = 1 - e^{-\lambda x}$ for all $x \in \R_+$. 
%%\EQ{
%%F(x) = P\set{X_{1}\le x} = \begin{cases}
%%1-e^{-\lambda x}, & x\ge 0  \\
%%0,  & \text{ else}.
%%\end{cases}
%%}
%%\end{defn}

For many proofs regarding Poisson processes, 
we partition the sample space with the disjoint events $\set{N_t = n}$ for $n \in \Z_+$. 
We need the following lemma that enables us to do that. 
\begin{lem} 
For any finite time $t > 0$, the number of points on the interval $(0,t]$ from a Poisson process is finite almost surely.
\end{lem}
\begin{proof} By strong law of large numbers, we have $\lim_{n \to \infty} \frac{S_{n}}{n} = \E[X_{1}] = \frac{1}{\lambda}$ almost surely. 
%Therefore, we have $S_n \to \infty$, a.s. This implies $P\set{N_t < \infty} =1$. To see this, 
Fix $t > 0$ and we define a sample space subset $M = \set{\omega \in \Omega: N(\omega, t) = \infty }$. 
For any $\omega \in M$, we have $S_{n}(\omega)\le t$ for all $n \in \N$. 
This implies $\lim\sup_n\frac{S_{n}}{n} = 0$  and $\omega \not\in \set{\lim_n \frac{S_{n}}{n} = \frac{1}{\lambda} }$. 
Hence, the probability measure for set $M$ is zero. 
\end{proof}

\subsection{Distribution functions}
\begin{lem} 
The following are true for the $n$th arrival instant $\tS_n$ of the Poisson arrival process $\tS:\Omega\to\R_+^\N$ with constant intensity density $\lambda$.
\begin{enumerate}[(a)]
\item The moment generating function is 
$
M_{\tS_n}(\theta) = \E[e^{\theta \tS_n} ] = \frac{\lambda^n}{(\lambda-\theta)^n}\SetIn{\theta < \lambda} + \infty\SetIn{\theta \ge \lambda}.
%\infty, & \theta \ge \lambda.
%\begin{cases}
%\frac{\lambda^n}{(\lambda-\theta)^n}, & \theta < \lambda \\ 
%\infty, & \theta \ge \lambda.
%\end{cases} 
$ 
\item The distribution function is  
$F_n(t) \triangleq P\set{\tS_n \le t} = 1 - e^{-\lambda t}\sum_{k=0}^{n-1}\frac{(\lambda t)^k}{k!}$.
\item The density function is Gamma distributed with parameters $n$ and $\lambda$. That is, $f_{n}(s) =\frac{\lambda (\lambda s)^{n-1}} {(n-1)!} e^{-\lambda s}$. 
\end{enumerate}
\end{lem}

% UNCOMMENT THIS %
%\begin{proof} 
%Since $S_n = \sum_{k=1}^nX_k$, where $X_k$ are \iid, the moment generating function $\mathbb{E} [ e^{\theta S_{n}} ]$ of $S_n$ is 
% \EQ{
%  \mathbb{E} [ e^{\theta S_{n}} ] = \left(\mathbb{E}[e^{\theta X_{1}}]\right)^{n}. 
% } 
%Since each $X_k$ is \iid exponential with rate $\lambda$, it is easy to see that moment generating function of inter-arrival time $X_1$ is 
% \EQ{
%  \mathbb{E} [ e^{\theta X_1} ] = 
%		\begin{cases}
%		\frac{\lambda}{(\lambda-\theta)}, & \theta < \lambda \\
%		\infty, & \theta \ge \lambda.
%		\end{cases} 
% } 
%\end{proof}
% UNCOMMENT THIS %
%\begin{proof} Notice that $X_i$'s are \iid and $S_1 = X_1$. In addition, we know that $S_n = X_n + S_{n-1}$. Since, $X_n$ is independent of $S_{n-1}$, we know that distribution of $S_n$ would be convolution of distribution of $S_{n-1}$ and $X_1$. Since $X_n$ and $S_1$ have identical distribution, we have $f_{n}=f_{n-1}*f_{1}$. The result follows from straightforward induction.
%\end{proof}

%Process $N_t$ is of real interest, and we can compute the distribution of $N_t$ for each $t$ from the distribution of $S_n$ in the following.
%\begin{figure}[hhhh]
%\center
	%\include{./Figures/Poisson}
 %% \caption{}\label{}
%\end{figure}
\begin{cor} 
Consider the counting process $N:\Omega\to\Z_+^{\R_+}$ associated with the Poisson arrival process $\tS:\Omega\to\R_+^\N$ having constant intensity density $\lambda$. 
The following are true.
\begin{enumerate}[(a)]
\item The relation between distribution of $n$th arrival instant and probability mass function for the counting process is given by  $F_n(t)= \sum_{j \ge n}P_j(t)$. 
\item For each $t\in\R_+$, the probability mass function $P_{N_t}\in\cM(\Z_+)$ for discrete random variable $N_t:\Omega\to\Z_+$ is given by $P_n(t) \triangleq P_{N_t}(n) = P\set{N_t=n)}= e^{-\lambda t}\frac{(\lambda t)^{n}}{n!}$. 
\item  The relation between distribution of $n$th arrival instant and the mean of the counting process is given by $\sum_{n \in \N}F_n(t) = \E N_t$. 
\item For each $t\in\R_+$, the mean  $\E[N_t] = \lambda t$, explaining the rate parameter $\lambda$ for the Poisson process. 
\end{enumerate}
\end{cor}
\begin{proof}
We observe the inverse relationship $\set{\tS_n \le t} = \set{N_t \ge n}$ for all $n \in \Z_+$ and $t\in\R_+$.  
\begin{enumerate}[(a)]
\item The result follows by taking the probability on both sides of the inverse relationship, to get  
$
F_n(t)  
= P\set{\tS_n \le t} 
= P\set{N_t \ge n}  
= \sum_{j \ge n}P\set{N_t = j}
= \sum_{j \ge n}P_j(t).
$
\item The result follows from the explicit from for the distribution of $\tS_n$ and recognizing that $P_n(t) = F_n(t) - F_{n+1}(t)$. 
\item The result from the following observation
$
\sum_{n \in \N}F_n(t) = \E\sum_{n \in \N}\SetIn{N_t \ge n} = \sum_{n \in \N}P\set{N_t \ge n} = \E N_t.
$
\item The result follows by summing the distribution function of the $n$th arrivals, to get 
$\E N_t = \sum_{n\in\N}F_n(t) = e^{-\lambda t}\sum_{n\in\N}\sum_{k \ge n}\frac{(\lambda t)^k}{k!} = 
\lambda t e^{-\lambda t}\sum_{k\in\N}\frac{(\lambda t)^{k-1}}{(k-1)!} = \lambda t$. 
\end{enumerate}
%\begin{eqnarray*}
%   P\set{N_t =n}&=&  P\set{S_{n}\le t, S_{n+1} >t)\\
%   &=&  \int^{t}_{0} P\left\set{ {S_{n+1}>t}|{S_{n}=s}\right}f_{n}(s)  ds\\
%   &\stackrel{(a)}{=}& \int^{t}_{0} P\set{X_{n+1}>t-s} f_{n}(s) ds\\
%   &=&  \int^{t}_{0}e^{-\lambda(t- s)} \frac{\lambda^{n}s^{n-1}}{(n-1)!}e^{-\lambda s}  ds\\
%   &=&\frac{e^{-\lambda t} (\lambda t)^{n}}{n !}.
%\end{eqnarray*}
% where (a) follows from the memoryless property of exponential distribution. %($P\set{X_{n+1}>s+t|X_{k+1}>t)=P\set{X_{k+1}>s)$).\\
\end{proof}

\begin{rem} 
A Poisson process is not a stationary process. 
That is, the finite dimensional distributions are not shift invariant. 
This is clear from looking at the first moment $\E N_t = \lambda t$, which is linearly increasing in time. 
\end{rem}
\end{document}